\section{Solver Specifications}
This section provides detailed descriptions of the five MILP solvers evaluated in our comparative study.

\subsection{GLPK (GNU Linear Programming Kit)}
GLPK is a free open-source linear programming solver developed as part of the GNU project \cite{glpk2024}. Version 4.65 was used in our experiments. GLPK uses the revised simplex method for linear programming and branch-and-cut algorithms for integer constraints. While GLPK provides reliable and deterministic solutions, it does not include the same heuristic or parallel features as commercial solvers, so it is slower for larger MILP instances. GLPK is widely used in academic research because it is accessible and performs well on small to medium-sized problems.

\subsection{HiGHS (High-performance Interior-point Solver)}
HiGHS (version 0.9) is an open-source linear and mixed-integer programming solver developed by the University of Edinburgh \cite{highs2023}. It uses simplex, interior-point, and active-set methods for LP problems, together with parallel branch-and-bound techniques for MILP. HiGHS is designed to be efficient and supports parallel execution. In practice, HiGHS is often faster than GLPK on many problem instances, especially when using multiple cores.

\subsection{SCIP (Solving Constraint Integer Programs)}
SCIP (version 7.0) is an open-source constraint integer programming framework developed by the Zuse Institute Berlin \cite{scip2023}. Beyond basic MILP, SCIP supports mixed-integer nonlinear programming and quadratically constrained problems. SCIP uses cutting planes, primal heuristics, and symmetry breaking. Among open-source solvers, SCIP often performs well and can be a strong option for researchers working with harder problem classes.

\subsection{Gurobi Optimizer}
Gurobi (version 10.0) is a commercial MILP solver developed by Gurobi Optimization \cite{gurobi2024}. It uses barrier methods (interior-point), parallel branch-and-cut, and presolve techniques. Gurobi also includes primal and dual heuristics and can distribute work across multiple processor cores. Its tuning and parameter options allow control over the solving process. It is widely used in industry and is a common benchmark for commercial solvers.

\subsection{CPLEX (IBM CPLEX Optimization Studio)}
CPLEX (version 12.10) is IBM's commercial optimization solver with decades of development \cite{cplex2024}. It uses primal and dual simplex, barrier methods with crossover, and branch-and-cut methods. CPLEX supports parallel processing and includes preprocessing and probing techniques. It is widely used in enterprise optimization and is one of the main commercial options for MILP solving.

\subsection{Experimental Configuration}
All solvers were executed with default parameters to ensure fair and reproducible comparison. This approach reflects typical usage patterns where practitioners rely on default settings without extensive parameter tuning. Each solver was allowed unlimited computation time (no time limits were imposed) to find optimal solutions. Experiments were conducted on an 11th Gen Intel Core i7-1165G7 processor operating at 2.80 GHz with 32 GB of system RAM, running Ubuntu Linux. This ensures all solvers had equal access to computational resources, with no artificial constraints or advantages.
